\subsubsection{Low-Temperature Intercept and Edge Weight Gap: Spectral Certificates for Maximum Weight Subkernel}\label{subsubsec:pom-low-temperature-maxedge-gap}
This section provides a general "low-temperature/large \(q\)" expansion for finite weighted directed graphs (or finite kernels): when edge weights are amplified by \(q\)-th powers, the logarithm of the spectral radius (pressure) is given by the exponential rate of the maximum edge weight and the Perron root of the maximum edge subgraph, with error exponentially decaying under the control of the edge weight gap \(\Delta\).
Furthermore, finitely many pressure samples at large \(q\) can directly extract verifiable certificates for the "ground state energy, ground state entropy, and gap" through finite differences.

\begin{definition}[Low-Temperature Moment Kernel and Pressure]\label{def:pom-low-temperature-weighted-adjacency}
Let \(G\) be a finite directed graph with vertex set \(V\) and edge set \(E\), and assign to each edge \(e\in E\) a weight \(a(e)>0\).
For real \(q\ge 0\), define the weighted adjacency matrix
\[
A(q)_{ij}:=\sum_{e:i\to j} a(e)^{q},
\qquad i,j\in V,
\]
and define the pressure function
\[
P(q):=\log\rho(A(q)).
\]
Write the maximum edge weight \(a_*:=\max_{e\in E}a(e)\), and let \(G_*\) be the subgraph generated by all edges satisfying \(a(e)=a_*\).
Its adjacency count matrix is denoted
\[
(A_*)_{ij}:=\#\{e:i\to j,\ a(e)=a_*\}.
\]
Let the second-largest edge weight
\[
a_{(2)}:=\max\{a(e):a(e)<a_*\},
\]
or by convention \(a_{(2)}=0\) if it does not exist.
Define the edge weight gap
\[
\Delta:=\log(a_*/a_{(2)})\in(0,\infty]\qquad(\log(a_*/0):=+\infty).
\]
\end{definition}

\begin{theorem}[Low-Temperature Three-Term Expansion and Exponential Error]\label{thm:pom-low-temperature-three-term-expansion}
Under the notation of Definition \ref{def:pom-low-temperature-weighted-adjacency}, assume \(\rho(A_*)>0\) (equivalent to \(G_*\) containing a directed cycle).
Then as \(q\to\infty\), we have
\[
\boxed{\;
P(q)=q\log a_*+\log\rho(A_*)+O(e^{-q\Delta})\; },
\]
where the error constant depends only on the graph size and number of edges.
More precisely, there exist constants \(C>0\) and \(q_0\ge 0\) such that for all \(q\ge q_0\),
\[
0\ \le\ P(q)-q\log a_*-\log\rho(A_*)\ \le\ C\,e^{-q\Delta}.
\]
\end{theorem}
\begin{proof}
Write \(A(q)\) as
\[
A(q)=a_*^{q}\bigl(A_*+E(q)\bigr),
\qquad
E(q)_{ij}:=\sum_{e:i\to j,\ a(e)<a_*}\left(\frac{a(e)}{a_*}\right)^q.
\]
By definition, \(E(q)\ge 0\) and there exists a constant \(C_0>0\) such that for any consistent matrix norm, \(\|E(q)\|\le C_0 e^{-q\Delta}\).
Since \(\rho\) is monotone on the non-negative cone, we have \(\rho(A_*+E(q))\ge \rho(A_*)\);
on the other hand, by \(\rho(B)\le \|B\|\) (for any submultiplicative norm) and \(\rho(A_*+E)\le \rho(A_*)+\rho(E)\), we obtain
\[
0\le \rho(A_*+E(q))-\rho(A_*)\le \rho(E(q))\le \|E(q)\|\le C_0 e^{-q\Delta}.
\]
Using \(\rho(A_*)>0\) and \(\log(1+x)=x+O(x^2)\) (as \(x\to 0\)), we get
\[
0\le \log\rho(A_*+E(q))-\log\rho(A_*)\le C\,e^{-q\Delta}
\]
for sufficiently large \(q\) (taking \(C\) to absorb constants).
From
\(
P(q)=q\log a_*+\log\rho(A_*+E(q))
\),
the conclusion follows.
\end{proof}

\begin{definition}[Zero-Temperature Gibbs-Markov Transition and Equilibrium State]\label{def:pom-zero-temperature-gibbs-markov}
Under the setup of Definition \ref{def:pom-low-temperature-weighted-adjacency}, further assume \(G\) is strongly connected, so that for each \(q\ge0\), the matrix \(A(q)\) is primitive.
Let \(r(q),\ell(q)\) be the left and right Perron vectors of \(A(q)\), normalized such that \(\ell(q)^\top r(q)=1\).
Define the Gibbs-Markov transition kernel
\[
\mathsf{P}_q(i\to j):=\frac{A(q)_{ij}\,r_j(q)}{\rho(A(q))\,r_i(q)},
\]
and define its stationary distribution (Perron equilibrium state)
\[
\pi_q(i):=\ell_i(q)\,r_i(q)\qquad(i\in V).
\]
\end{definition}

\begin{theorem}[Zero-Temperature Gibbs State Perron Vector Rigidity: Exponential Concentration to Maximum Weight Subkernel]\label{thm:pom-zero-temperature-gibbs-rigidity}
Adopt the notation of Definitions \ref{def:pom-low-temperature-weighted-adjacency} and \ref{def:pom-zero-temperature-gibbs-markov}.
Let \(G_*\) be the maximum edge weight subgraph, and let \(V_*\subseteq V\) be its vertex set (equivalently: there exists \(j\) such that \((A_*)_{ij}+(A_*)_{ji}>0\)).
Further assume that \(A_*\) viewed as acting on \(\RR^{V_*}\) is a primitive matrix.
Then there exist constants \(C>0\) and \(q_0\) such that for all \(q\ge q_0\),
\[
\boxed{\;
\sum_{i\notin V_*}\pi_q(i)\ \le\ C\,e^{-q\Delta}\; },
\qquad
\boxed{\;
\max_{i\notin V_*}\frac{r_i(q)}{\max_{v\in V_*}r_v(q)}\ \le\ C\,e^{-q\Delta}\; },
\qquad
\boxed{\;
\max_{i\notin V_*}\frac{\ell_i(q)}{\max_{v\in V_*}\ell_v(q)}\ \le\ C\,e^{-q\Delta}\; }.
\]
Therefore, as \(q\to\infty\), the zero-temperature equilibrium state concentrates exponentially fast to the maximum weight subkernel \(G_*\), with the concentration rate precisely controlled by the edge weight gap \(\Delta\).
\end{theorem}
\begin{proof}
Write \(A(q)=a_*^q(A_*+E(q))\), where \(E(q)\ge0\) and \(\|E(q)\|\le C_0e^{-q\Delta}\) (as given in the proof of Theorem \ref{thm:pom-low-temperature-three-term-expansion}).
Since \(A_*\) (on \(V_*\)) is primitive, its Perron projection is spectrally separated from the remaining spectral points.
Applying standard analytic perturbation estimates for non-negative matrices to the Perron root and Perron vectors, we find that \(r(q),\ell(q)\) converge on \(V_*\) to the Perron vectors of \(A_*\), while the leakage to the complement space is controlled by the perturbation magnitude \(O(e^{-q\Delta})\).
Substituting this estimate into \(\pi_q=\ell\cdot r\) yields the conclusion.
\end{proof}

\begin{corollary}[Zero-Temperature Limit Equilibrium State and TV Exponential Convergence]\label{cor:pom-zero-temperature-gibbs-limit}
Under the conditions of Theorem \ref{thm:pom-zero-temperature-gibbs-rigidity}, there exists a limit stationary distribution \(\pi_\infty\) supported on \(V_*\), and
\[
D_{\mathrm{TV}}(\pi_q,\pi_\infty)\le C\,e^{-q\Delta}
\]
holds for sufficiently large \(q\), where \(C\) is of the same order as in Theorem \ref{thm:pom-zero-temperature-gibbs-rigidity}.
\end{corollary}
\begin{proof}
By Theorem \ref{thm:pom-zero-temperature-gibbs-rigidity}, \(\{\pi_q\}\) is a Cauchy sequence in the finite-dimensional simplex with exponentially decaying mass leakage, so a limit \(\pi_\infty\) exists with \(\mathrm{supp}(\pi_\infty)\subseteq V_*\).
The TV distance bound follows from the mass leakage and exponential convergence on \(V_*\).
\end{proof}

\begin{theorem}[Difference Certificates: Extracting Intercept and Gap from Finitely Many Pressure Values at Large \texorpdfstring{$q$}{q}]\label{thm:pom-low-temperature-difference-certificates}
Under the setup of Theorem \ref{thm:pom-low-temperature-three-term-expansion}, for integer \(q\ge 1\) define the difference quantities
\[
\widehat\alpha(q):=P(q+1)-P(q),\qquad
\widehat h(q):=P(q)-q\,\widehat\alpha(q),\qquad
\Delta^2P(q):=P(q+1)-2P(q)+P(q-1).
\]
Then there exist constants \(C>0\) and \(q_0\) such that for all integers \(q\ge q_0\),
\[
\boxed{\;
\bigl|\widehat\alpha(q)-\log a_*\bigr|\le C\,e^{-q\Delta}\; },
\qquad
\boxed{\;
\bigl|\widehat h(q)-\log\rho(A_*)\bigr|\le C\,q\,e^{-q\Delta}\; }.
\]
Moreover,
\[
\boxed{\;
\Delta^2P(q)=O(e^{-q\Delta})\qquad(q\to\infty)\; }.
\]
Furthermore, if there exists a constant \(c\neq 0\) such that \(\Delta^2P(q)\sim c\,e^{-q\Delta}\), then the gap satisfies the verifiable limit formula
\[
\boxed{\;
\Delta=-\lim_{q\to\infty}\log\left|\frac{\Delta^2P(q+1)}{\Delta^2P(q)}\right|\; }.
\]
\end{theorem}
\begin{proof}
By Theorem \ref{thm:pom-low-temperature-three-term-expansion}, write
\(
P(q)=q\log a_*+\log\rho(A_*)+R(q)
\)
where \(R(q)=O(e^{-q\Delta})\).
Then
\[
\widehat\alpha(q)-\log a_*=R(q+1)-R(q)=O(e^{-q\Delta}),
\]
and
\[
\widehat h(q)-\log\rho(A_*)
=R(q)-q\bigl(R(q+1)-R(q)\bigr)
=O(e^{-q\Delta})+q\,O(e^{-q\Delta})
=O(qe^{-q\Delta}).
\]
Similarly,
\(
\Delta^2P(q)=R(q+1)-2R(q)+R(q-1)=O(e^{-q\Delta})
\).
If further \(\Delta^2P(q)\sim c e^{-q\Delta}\), then
\(
\Delta^2P(q+1)/\Delta^2P(q)\to e^{-\Delta}
\),
taking logarithm yields the limit formula.
\end{proof}

\begin{remark}[Tropical-Perron Binary Identification: Maximum Cycle Average and Critical Graph Intercept]\label{rem:pom-tropical-perron-binary-identification}
Generally, let \(w(e):=\log a(e)\). For each directed cycle \(\gamma=e_1\cdots e_n\), define its logarithmic average weight
\[
\overline{w}(\gamma):=\frac{1}{n}\sum_{t=1}^{n}w(e_t),
\qquad
\alpha_{\max}:=\max_{\gamma}\ \overline{w}(\gamma)
\]
(max-plus eigenvalue/maximum cycle average).
Let the critical graph \(G_{\mathrm{crit}}\) be the subgraph generated by edges contained in cycles achieving \(\alpha_{\max}\), and let \(A_{\mathrm{crit}}\) be its adjacency count matrix.
Then for the pressure \(P(q)=\log\rho(A(q))\), we have the limit identification
\[
\alpha_\ast:=\lim_{q\to\infty}\frac{P(q)}{q}=\alpha_{\max},
\qquad
h_\ast:=\lim_{q\to\infty}\bigl(P(q)-q\alpha_\ast\bigr)=\log\rho(A_{\mathrm{crit}}),
\]
and under the maximum edge gap regime of this section (\(\rho(A_*)>0\)), we have \(\alpha_{\max}=\log a_*\), and the spectral radius of \(A_{\mathrm{crit}}\) coincides with that of \(A_*\), thus recovering the two-term expansion form of Theorem \ref{thm:pom-low-temperature-three-term-expansion}.
\end{remark}

\begin{remark}[Exponential Precision Approximation of Zero-Temperature Intercept by Truncated Cyclotomic Atoms (Coupling Regime)]\label{rem:pom-zero-temperature-intercept-cyclotomic-truncation}
The zero-temperature intercept \(h_\ast\) can be understood as the topological entropy of a certain "maximum phase/critical subkernel" (Theorem \ref{thm:pom-low-temperature-three-term-expansion} and Corollary \ref{cor:pom-max-fiber-subshift-spectral-entropy}).
On the other hand, the constant-layer cyclotomic atom sequence \(\mathbf{C}(M)=(\mathcal{C}_r(M))_{r\ge2}\) constitutes complete coordinates for the normalized non-Perron spectrum (Theorem \ref{thm:finite-part-cyclotomic-atoms-complete-coordinate}).
\par
When both of the following regimes are simultaneously satisfied:
\begin{enumerate}
  \item \textbf{(Zero-temperature gap regime)} There exists \(\Delta>0\) such that the low-temperature/zero-temperature error terms satisfy exponential closure (e.g., the \(e^{-q\Delta}\) scale of Theorems \ref{thm:pom-low-temperature-three-term-expansion} and \ref{thm:pom-zero-temperature-gibbs-rigidity});
  \item \textbf{(Constant-layer coupling regime)} The critical subkernel (hence \(h_\ast\)) can be read out through an explicit algebraic function \(R_\infty\) from the constant-layer spectral coordinates \(\mathbf{C}\) of some finite-dimensional matrix,
\end{enumerate}
then \(h_\ast\) can in principle be approximated to exponential precision by finitely many atom truncations: there exists an explicitly computable function \(R_Q\) (depending only on \(\mathcal{C}_2,\dots,\mathcal{C}_Q\) and the gap parameter \(\Delta\)) such that
\[
\bigl|h_\ast-R_Q(\mathcal{C}_2,\dots,\mathcal{C}_Q)\bigr|\ \le\ K\,e^{-Q\Delta},
\]
where the constant \(K\) depends only on the finite state kernel size and the chosen readout implementation.
This conclusion can be viewed as a verifiable connectivity law for "zero-temperature structure \(\leftrightarrow\) constant-layer spectral coordinates": under the gap regime, finite-depth constant-layer atoms suffice to lock down the exponential precision of \(h_\ast\).
\end{remark}

